% Generated draft: requires mathematical and source review.
\documentclass[11pt]{article}
\usepackage[margin=25mm]{geometry}
\usepackage{fontspec}
\setmainfont{Latin Modern Roman}
\usepackage{amsmath,amssymb,mathtools,unicode-math}
\setmathfont{Latin Modern Math}
\usepackage{enumitem,needspace,xurl,hyperref}
\setlength{\parindent}{0pt}
\setlength{\parskip}{5pt}
\setlength{\emergencystretch}{3em}
\newcommand{\N}{\mathbb{N}}
\newcommand{\Z}{\mathbb{Z}}
\newcommand{\Q}{\mathbb{Q}}
\newcommand{\R}{\mathbb{R}}
\newcommand{\C}{\mathbb{C}}
\newcommand{\F}{\mathbb{F}}
\newcommand{\A}{\mathbb{A}}
\newcommand{\PP}{\mathbb{P}}
\newcommand{\E}{\mathbb{E}}
\newcommand{\eps}{\varepsilon}
\newcommand{\dd}{\,\mathrm{d}}
\begin{document}
\section*{UnsolvedMath ID: 30006195}
\textbf{Source problem code:} catalogue entry
\subsection{Definitions and mathematical statement}
A Kähler group is a group isomorphic to \(\pi_1(X,x)\) for a connected compact Kähler manifold \(X\). A group \(G\) is residually finite if, for every \(g\in G\setminus\{1\}\), there exist a finite group \(Q\) and a homomorphism \(f:G\to Q\) with \(f(g)\ne 1\). Its profinite completion is
\[
\widehat G=\varprojlim_{\substack{N\trianglelefteq G\\[G:N]<\infty}}G/N,
\]
with its inverse-limit topology. Let \(\mathcal K\) denote the class of Kähler groups and \(\mathcal K_{\mathrm{rf}}\) its subclass of residually finite groups. For either comparison class \(\mathcal C\), define
\[
\operatorname{Gen}_{\mathcal C}(G)
=\bigl\{[H]_{\cong}:H\in\mathcal C,\ \widehat H\cong_{\mathrm{top}}\widehat G\bigr\},
\]
where \([H]_{\cong}\) is the abstract group isomorphism class of \(H\), and \(\cong_{\mathrm{top}}\) means isomorphism of topological groups.

The first supplied question asks whether
\[
\forall G\in\mathcal K_{\mathrm{rf}},\qquad
\bigl|\operatorname{Gen}_{\mathcal C}(G)\bigr|<\infty.
\]
The supplied record calls this the finiteness of the \(\mathcal K\)-profinite genus, but does not supply its definition. In particular, it does not establish whether the comparison class is \(\mathcal C=\mathcal K_{\mathrm{rf}}\) or \(\mathcal C=\mathcal K\). These are distinct candidate formulations, and the intended choice requires source clarification.

A group \(S\) is simple if \(S\ne\{1\}\) and its only normal subgroups are \(\{1\}\) and \(S\). The second supplied question, read literally, is
\[
\exists S\in\mathcal K\quad S\text{ is simple}.
\]
The record does not explicitly require \(S\) to be infinite. If infinite simplicity is intended, the formulation instead is
\[
\exists S\in\mathcal K\quad
\bigl(|S|=\infty\ \text{and}\ S\text{ is simple}\bigr).
\]
This additional condition is unresolved here. Residual finiteness is imposed on \(G\) in the first question; the supplied wording imposes no such condition on \(S\) in the second.

The background also mentions the related question
\[
\exists H\in\mathcal K\quad H\text{ is not residually finite}.
\]
The supplied curation attributes a positive answer to Toledo. This is contextual information, not an additional unresolved target or an independently verified result in this draft.
\subsection{Short English statement}
The record asks whether every residually finite Kähler group has only finitely many Kähler groups with the same profinite completion, counted up to group isomorphism, and whether a simple Kähler group exists. The comparison class for the genus and any intended infinitude requirement for the simple group need clarification. Supplied dated curation attributes the related existence of a non-residually-finite Kähler group to Toledo.
\subsection{Sources}
\begin{itemize}
\item Local registry: \texttt{data/unsolvedmath/problems.json}, permanent UnsolvedMath ID \(30006195\), exact source code \texttt{catalogue entry}; title: \emph{Profinite Genera and Simple Kähler Groups}. The supplied record identifies the collection as \emph{Oberwolfach Reports Open Problems}, with proposed year 2025 and no named proposer.
\item Supplied report citation: \emph{Geometric Structures in Group Theory} (2025), Oberwolfach Reports, pp.\ 275--338, DOI \texttt{10.4171/owr/2025/6}. \url{https://doi.org/10.4171/owr/2025/6}.
\item Supplied literature source, labelled ``OWR 2025/6 Kahler-group report'': \url{https://ora.ox.ac.uk/objects/uuid%3A3212138f-4ec5-43d4-9af3-49789a6dce51/files/r6395w9580}.
\item Supplied literature source, labelled ``Hughes et al., profinite rigidity of Kahler groups'': \url{https://arxiv.org/abs/2501.13761}.
\item Supplied literature source, labelled ``Toledo, non-residually-finite Kahler groups'': \url{https://library.slmath.org/books/Book37/files/toledo.pdf}. The extracted report attributes the positive answer to Toledo, reference [16], and describes the simple-group question as a variation of Serre's question about non-residually-finite Kähler groups.
\item Literature assessment dated exactly \texttt{2026-08-22}: supplied dated curation, not independently verified here. It labels the record \texttt{partially\_solved}, reports that genus finiteness and existence of a simple Kähler group remain open in the report, and attributes a positive answer to the related non-residual-finiteness question to Toledo. Its background expressly requests independent verification of time-sensitive status and cited claims.
\item The registry labels the statement \texttt{corrected\_verified} and says it was reconstructed and checked against the official report. Those are supplied provenance claims; this draft does not independently verify them or claim to have visited any source.
\end{itemize}
\end{document}
